% ---------------------------------------------------------------------------
% Standing Water — an illustrated collection, published 1898.
%
% Typographic premise: an American trade octavo of the late 1890s, illustrated
% with mounted photographic plates. NOT an antiquarian private press — the
% Necronomicon pipeline this is grown from is set in the Fell types, which are
% 1680s and Oxford, and would be two centuries and an ocean wrong here.
%
% Body face: Old Standard (OFL), a digitisation of the "modern" face that set
% most English-language books between about 1860 and 1930. It has no true
% small capitals and no bold in the period sense, which suits: emphasis in a
% book of 1898 is italic, and headings are letterspaced capitals.
% ---------------------------------------------------------------------------

\usepackage{fontspec}
\usepackage[
  paperwidth=5.5in, paperheight=8.25in,
  inner=0.78in, outer=0.95in, top=0.80in, bottom=1.05in,
  headsep=14pt, footskip=26pt
]{geometry}
\usepackage{lettrine}
\usepackage{fancyhdr}
\usepackage{microtype}
\usepackage{ragged2e}
\usepackage{graphicx}
\usepackage{needspace}
\usepackage{eso-pic}
\graphicspath{{plates/}}

% Relative, so the repo can move. The build runs xelatex with cwd = render/.
\def\fontdir{fonts/}

% The body face is chosen by tex/face.tex, which scripts/set_face.py writes
% from tex/faces/<name>.tex. Without one the default below applies. Faces
% differ in more than shape: some have real small capitals and some do not,
% so each face file also says which \caps to use.
\InputIfFileExists{tex/face}{}{%
  \setmainfont{OldStandard-Regular.ttf}[
    Path = \fontdir , ItalicFont = OldStandard-Italic.ttf ,
    BoldFont = OldStandard-Bold.ttf , Ligatures = TeX ]}

% ------------------------------------------------------------------ measure
\linespread{1.16}
\setlength{\parindent}{1.15em}
\frenchspacing
\raggedbottom
\clubpenalty=8000
\widowpenalty=8000

% ------------------------------------------------------------ letterspacing
% microtype's \textls needs pdfTeX or LuaTeX; under XeTeX the letterspacing
% has to come from the font engine, which fontspec exposes as LetterSpace.
\newcommand{\ls}[2]{{\addfontfeature{LetterSpace=#1}#2}}

% Small capitals. Some of the candidate faces have a real cut (IM Fell, Sorts
% Mill Goudy) and some have none at all (Old Standard, Libre Caslon), so both
% are defined here and each face file says which one it gets. The synthesised
% version is uppercase at 82% of the current size, letterspaced — what a
% photo-set fake has always done, and adequate at the sizes used here.
\newlength{\capsize}
\makeatletter
% \f@size is a bare number, not a dimen, so it cannot be scaled in place:
% 0.82\f@size expands to "0.8210.95" and TeX reports an illegal unit. It has
% to go through a length register first.
\newcommand{\fakecaps}[1]{%
  {\addfontfeature{LetterSpace=6.0}%
   \setlength{\capsize}{\f@size pt}%
   \setlength{\capsize}{0.82\capsize}%
   \fontsize{\capsize}{\f@baselineskip}\selectfont
   \MakeUppercase{#1}}}
\newcommand{\realcaps}[1]{{\addfontfeature{LetterSpace=4.0}\scshape #1}}
\makeatother
% Faces with a real small-capitals cut override this in tex/faces/<name>.tex.
\let\caps\fakecaps

% ------------------------------------------------------------------ story head
% Each story opens on a fresh recto, sunk about an inch, title in letterspaced
% capitals over a short rule. \storyend clears to the next recto and leaves the
% verso blank, which is how a collection of this period sets its stories.
\newcommand{\storyhead}[1]{%
  \cleardoublepage
  \thispagestyle{storyopen}%
  \null\vspace*{0.55in}%
  \begin{center}
    {\large\ls{9.0}{\MakeUppercase{#1}}}\par\vspace{0.55\baselineskip}
    \rule{0.16\textwidth}{0.4pt}
  \end{center}
  \vspace{1.05\baselineskip}%
  \markboth{Standing Water}{#1}%
  \pagestyle{fancy}%
  \noindent\ignorespaces}

\newcommand{\storyend}{\cleardoublepage}

% A scene break: a little air and a short centred rule. Not three asterisks —
% that is a modern typescript convention, not a printed one.
\makeatletter
\newcommand{\scenebreak}{%
  \par\addvspace{0.9\baselineskip}%
  \needspace{3\baselineskip}%
  \noindent\hfill\rule{0.10\textwidth}{0.4pt}\hfill\null\par
  \addvspace{0.9\baselineskip}%
  \@afterindentfalse\@afterheading}
\makeatother

% ------------------------------------------------------------------ drop cap
\setcounter{DefaultLines}{2}
\renewcommand{\LettrineFontHook}{\normalfont}
% lettrine sets its run-in text in small capitals by default; the text face
% has none, so the run-in is left plain and \caps does the work instead.
\renewcommand{\LettrineTextFont}{\normalfont}
\setlength{\DefaultFindent}{2pt}
\setlength{\DefaultNindent}{0pt}

% ---------------------------------------------------------- blank versos
% \cleardoublepage inserts a blank verso but does not clear its page style, so
% the blank inherits the running head of the story just ended — which then
% appears after that story is over. Blanks must be truly blank.
\makeatletter
\renewcommand*{\cleardoublepage}{%
  \clearpage
  \if@twoside
    \ifodd\c@page\else
      \null\thispagestyle{empty}\newpage
    \fi
  \fi}
\makeatother

% ------------------------------------------------------------- running heads
% Folio in the outer corner of the head, running title centred: the standard
% arrangement for an illustrated trade octavo of the nineties. Nothing at the
% foot, so a plate leaf reads as blank paper below the caption.
\pagestyle{fancy}
\fancyhf{}
\renewcommand{\headrulewidth}{0pt}
\fancyhead[LE]{{\footnotesize\thepage}}
\fancyhead[RO]{{\footnotesize\thepage}}
\fancyhead[CE]{{\footnotesize\ls{5.0}{\MakeUppercase{\leftmark}}}}
\fancyhead[CO]{{\footnotesize\ls{5.0}{\MakeUppercase{\rightmark}}}}

\fancypagestyle{storyopen}{%
  \fancyhf{}\renewcommand{\headrulewidth}{0pt}%
  \fancyfoot[C]{{\footnotesize\thepage}}}

% -------------------------------------------------------------- front matter
\newcommand{\fmhead}[1]{%
  \cleardoublepage
  \thispagestyle{storyopen}%
  \null\vspace*{0.45in}%
  \begin{center}
    {\ls{8.0}{\MakeUppercase{#1}}}\par\vspace{0.45\baselineskip}
    \rule{0.14\textwidth}{0.4pt}
  \end{center}
  \vspace{1.0\baselineskip}%
  \markboth{#1}{#1}}

% ------------------------------------------------------------------- plates
% A mounted plate: its own leaf, no folio, no running head, printed one side
% only — so the blank verso that follows is correct rather than wasteful. The
% caption sits below in small type, the plate numeral run in at its head.
%
% The ageing pass finds these pages by detecting colour: the text is black on
% white, a plate is a warm-toned photograph. Nothing here records page numbers.
%
% The height cap is load-bearing, not taste. The plate, its caption and the
% fill above and below have to fit one 	extheight. At 0.80 they did not:
% the box spilled onto the following recto, and the story — which opens with
% \cleardoublepage — was pushed a further leaf on and no longer faced its
% plate. If the plate is ever enlarged again, check that a story still opens
% on the recto opposite it.
%
%   \plate{file}{numeral}{caption}
%
% A plate is tipped in FACING the story it belongs to, which means it prints on
% a verso: land on a recto, leave that recto blank (the back of the inserted
% leaf), print the plate on the verso, and the story then opens on the recto
% opposite. Placing it on a recto instead — the obvious reading of "its own
% leaf" — puts a blank page between the photograph and the story it is of.
\newcommand{\plate}[3]{%
  \cleardoublepage
  \null\thispagestyle{empty}\clearpage
  \thispagestyle{empty}%
  \null\vfill
  \begin{center}
    \includegraphics[width=1.30\textwidth,height=0.72\textheight,keepaspectratio]{#1}%
    \platecaption{#2}{#3}%
  \end{center}
  \vfill\null
  \clearpage}

\newcommand{\platecaption}[2]{%
  \par\vspace{1.1\baselineskip}
  \begin{minipage}{0.88\textwidth}\footnotesize\centering
    \ls{4.0}{\MakeUppercase{Plate #1}}\par\vspace{0.35\baselineskip}
    \RaggedRight #2
  \end{minipage}}

% The frontispiece is the exception: it faces the title page, so it prints on
% the verso immediately after the half-title, with no inserted blank and no
% caption — an empty \plate would otherwise set the word PLATE above nothing.
\newcommand{\frontispieceplate}[1]{%
  \clearpage
  \thispagestyle{empty}%
  \null\vfill
  \begin{center}
    \includegraphics[width=0.92\textwidth,height=0.78\textheight,keepaspectratio]{#1}%
  \end{center}
  \vfill\null
  \clearpage}

% ------------------------------------------------------- full-bleed plate
% A leaf that IS the photograph: printed to the trim, no margin, no folio, no
% caption. The source is already a finished leaf — a print on its own mount,
% photographed — so setting it inside the text measure would paste one sheet
% of paper onto another, and the join shows as a rectangle.
%
% eso-pic paints it into the page background, which is the only reliable way
% to reach outside the type area under twoside, where the margin an \hspace
% would have to cancel differs between recto and verso.
%
%   \platefull{file}
\newcommand{\platefull}[1]{%
  \cleardoublepage
  \null\thispagestyle{empty}\clearpage   % blank recto: the back of the leaf
  \thispagestyle{empty}%
  \AddToShipoutPictureBG*{%
    \put(0,0){\includegraphics[width=\paperwidth,height=\paperheight]{#1}}}%
  \null\clearpage}

% The boards, as the first leaf, with no blank in front of them.
\newcommand{\coverleaf}[1]{%
  \thispagestyle{empty}%
  \AddToShipoutPictureBG*{%
    \put(0,0){\includegraphics[width=\paperwidth,height=\paperheight]{#1}}}%
  \null\clearpage
  \null\thispagestyle{empty}\clearpage}   % the inside of the board
